\documentclass[10pt,a4paper]{amsart}
\usepackage[margin=19mm]{geometry}
\usepackage{amsmath,amssymb,mathtools}
\usepackage[scr=rsfs,cal=euler]{mathalfa}
\usepackage{xfrac}
\usepackage[unicode,hidelinks,bookmarks=false]{hyperref}
\newcommand{\ie}{i.e.\ }
\newcommand{\cf}{cf.\ }
\newcommand{\resp}{respectively\ }
\newcommand{\Subsection}{\subsection}
\providecommand{\nobreakdash}{\nobreak\hbox{-}}
\setlength{\emergencystretch}{2em}
\multlinegap0pt
\usepackage{mathtools}
\usepackage{amssymb}
\usepackage{enumerate}
\newtheorem{theorem}{Theorem}
\newtheorem{proposition}{Proposition}
\newcommand{\C}{\mathbb{C}}
\newcommand{\R}{\mathbb{R}}
\DeclareMathOperator{\Res}{Res}
\newcommand{\difrac}{\displaystyle \frac}
\newcommand{\diint}{\displaystyle \int}
\newcommand{\dilim}{\displaystyle \lim}
\newcommand{\disum}{\displaystyle \sum}
\title{Real exponential sums over primes and prime gaps}
\author{Luan Alberto Ferreira}
\date{Source: arXiv:2307.08725v4 (CC BY 4.0)}
\begin{document}
\maketitle

\noindent\emph{Source and licence.} Real exponential sums over primes and prime gaps. Version arXiv:2307.08725v4 (CC BY 4.0), distributed by arXiv under the Creative Commons Attribution 4.0 International licence, http://creativecommons.org/licenses/by/4.0/, as recorded in the arXiv licence metadata for this exact version and shown on the abstract page https://arxiv.org/abs/2307.08725v4. Full licence notice: \texttt{licenses/arxiv-2307.08725-CC-BY-4.0.txt}.\par\smallskip

\noindent\textit{Editorial scope.} Counted unit: the article's proposition proposition (source0.tex lines 468--472). The article's complete original proof of it is delivered with it (source0.tex lines 474--514, one \texttt{proof} environment of 41 lines). No mathematics is written, repaired or re-derived: the statement and the proof are the article's own lines, in the article's own notation, and no retained line is substituted or re-flowed.  Retained uncounted as the source-local context the proof needs: lines 444--448; lines 454--462.  Nothing was omitted from any retained span, so the package carries no omission record.  No retained line contains a citation, so the wrapper carries no bibliography and no bibliography entry.  No retained line contains a figure environment, an image-inclusion command or drawing code, so no image is delivered and the package declares no asset dependency.  Editorial formatting only, in addition: an AMS \texttt{amsart} preamble that restates the article's own declarations and macros used by the retained lines, namely the environments \texttt{theorem}, \texttt{proposition} on separate counters, together with the standard packages those lines need.  The retained environments print in this document as Theorem 1, Proposition 1 in order of appearance, whereas the article prints the same items with the numbering of its own full document.  The complete native source of the article as distributed in the arXiv e-print for this exact version is bundled unchanged as \texttt{sources/144401/source0.tex}.
\begin{theorem} \label{wellknowntheorem}

The function $$\Phi(s + 1) - \difrac{1}{s} = \disum_p \difrac{\log(p)}{p^{s + 1}} - \difrac{1}{s},$$ initially defined on $\Re(s) > 0$, extends to a meromorphic function on $\Re(s) > -\difrac{1}{2}$ that is analytic on $\Re(s) \geq 0$.

\end{theorem}

\begin{theorem} \label{wellknowntheorem2}

There exists a constant $b > 0$ such that the Riemann zeta function has no zeros in the region $$\left\{s = \sigma + it \in \C; \ \sigma > 1 - \difrac{b}{\log(2+|t|)}\right\}.$$

\noindent Consequently, there exists a constant $a > 0$ such that $\Omega(s)$ is analytic on the open set $$\mathcal{U} = \left\{s = \sigma + it \in \C; \ \sigma > - \difrac{a}{\log(2+|t|)}\right\}.$$

\noindent Furthermore, within the open neighborhood of the imaginary axis defined by $$\mathcal{V} = \left\{s = \sigma + it \in \C; \ |\sigma| < \difrac{a}{\log(2+|t|)}\right\},$$ the function $\Omega(s)$ satisfies the growth estimate $$\Omega(s) = O(\log(2 + |t|)).$$ 

\end{theorem}

\begin{proposition}

The function $\tau(s) - \difrac{1}{s}$, initially defined for $\Re(s) > 0$, extends analytically to $\Re(s) \geq 0$.

\end{proposition}

\begin{proof} As mentioned above, calculating the Mellin transform of the function $T(s)$ for $\Re(s) > 0$, we obtain \begin{equation*} \begin{split} \mathscr{M}(z) & = \mathscr{M}\left(-\disum_p \difrac{(1 - \lambda) \log(p)}{p^{3 - 2\lambda} \exp(sp^{1 - \lambda})}\right) \\ & = -\disum_p \difrac{(1 - \lambda) \log(p)}{p^{3 - 2\lambda}} \mathscr{M}[\exp(-sp^{1 - \lambda})] \\ & = -\disum_p \difrac{(1 - \lambda) \log(p)}{p^{3 - 2\lambda}} \cdot \difrac{\Gamma(z)}{p^{(1 - \lambda)z}} \\ & = -(1 - \lambda) \Gamma(z) \disum_p \difrac{\log(p)}{p^{1 + (1 - \lambda)(z + 2)}} \\ & = -(1 - \lambda) \Gamma(z) \Phi(1 + (1 - \lambda)(z + 2)). \end{split} \end{equation*}

Note that this function is meromorphic on $\Re(z) > -2$, and has two simple poles: one at $z = 0$ and another at $z = -1$; not because of the function $\Phi(1 + (1 - \lambda)(z + 2))$, which is analytic on $\Re(z) > -2$, but because of the function $\Gamma(z)$ with its two simple poles at $z = 0$ and $z = -1$. Using the information from theorem \ref{wellknowntheorem}, we can write \begin{equation*} \begin{split} \mathscr{M}(z) & = -(1 - \lambda) \Gamma(z) \left[\difrac{1}{(1 - \lambda)(z + 2)} + \Omega((1 - \lambda)(z + 2))\right] \\ & = -\difrac{\Gamma(z)}{z + 2} - (1 - \lambda) \Gamma(z) \Omega((1 - \lambda)(z + 2)). \end{split} \end{equation*}

Now, for $\Re(z) > -2$, notice that this function has exponential decay along vertical lines, because $\Gamma(z)$ does, and that the function $\Omega((1 - \lambda)(z + 2))$ doesn't change that . Therefore, for $\Re(z) > -2$, we can get back the function $T(s)$ by computing the inverse Mellin transform of $\mathscr{M}(z)$, and write \begin{equation*} \begin{split} T(s) & = \mathscr{M}^{-1}(\mathscr{M}(z)) \\ & = \mathscr{M}^{-1}\left(-\difrac{\Gamma(z)}{z + 2} - (1 - \lambda) \Gamma(z) \Omega((1 - \lambda)(z + 2))\right) \\ & = -\mathscr{M}^{-1}\left(\difrac{\Gamma(z)}{z + 2}\right) - (1 - \lambda) \mathscr{M}^{-1}[\Gamma(z) \Omega((1 - \lambda)(z + 2))] \\ & = -\difrac{s^2 \cdot \Gamma(0,s) + \exp(-s) \cdot (1-s)}{2} - (1 - \lambda) \mathscr{M}^{-1}[\Gamma(z) \Omega((1 - \lambda)(z + 2))] \end{split} \end{equation*}

\noindent for all $s \in \C$ such that $\Re(s) > 0$. Differentiating three times, one gets \begin{equation*} \begin{split} \tau(s) & = T'''(s) \\ & = \difrac{d^3}{ds^3} \left(-\difrac{s^2 \cdot \Gamma(0,s) + \exp(-s) \cdot (1-s)}{2}\right) - (1 - \lambda) \difrac{d^3}{ds^3} \left(\mathscr{M}^{-1}[\Gamma(z) \Omega((1 - \lambda)(z + 2))]\right) \\ & = \difrac{\exp(-s)}{s} - (1 - \lambda) \difrac{d^3}{ds^3} \left(\mathscr{M}^{-1}[\Gamma(z) \Omega((1 - \lambda)(z + 2))]\right). \end{split} \end{equation*}

Subtracting $\difrac{1}{s}$ on either side of this equation, we get $$\tau(s) - \difrac{1}{s} = \difrac{\exp(-s) - 1}{s} - (1 - \lambda) \difrac{d^3}{ds^3} \left(\mathscr{M}^{-1}[\Gamma(z) \Omega((1 - \lambda)(z + 2))]\right).$$

But we have already seen that the function $$s \mapsto \difrac{\exp(-s) - 1}{s}$$ extends analytically to $\C$. The key observation is that the function $$\mathscr{M}^{-1}[\Gamma(z) \Omega((1 - \lambda)(z + 2))] = \difrac{1}{2\pi i} \diint_{c - i\infty}^{c + i\infty} s^{-z} \cdot \Gamma(z) \cdot \Omega((1-\lambda)(z+2)) dz,$$ initially defined and analytic on $\Re(s) > 0$, can be analytically extended to $\Re(s) \geq 0$ via a deformation of the contour $\mathcal{C}$ of the integral for the points that are on (and close to the left of) the line $\Re(z) = -2$. Consequently, the function $$\difrac{d^3}{ds^3} \left(\mathscr{M}^{-1}[\Gamma(z) \Omega((1 - \lambda)(z + 2))]\right),$$ being the third derivative of an analytic function in $\Re(s) \geq 0$, will itself be analytic in $\Re(s) \geq 0$, thus concluding the proof of this proposition.

With this in mind, by making the change of variables $z + 2 \leadsto z$ and using the classic property of the $\Gamma$ function, we can write \begin{equation*} \begin{split} \mathscr{M}^{-1}[\Gamma(z) \Omega((1 - \lambda)(z + 2))] & = \difrac{1}{2\pi i} \diint_{c - i\infty}^{c + i\infty} s^{-z} \cdot \Gamma(z) \cdot \Omega((1-\lambda)(z+2)) dz \\ & = \difrac{1}{2\pi i} \diint_{c - i\infty}^{c + i\infty} s^{-(z - 2)} \cdot \Gamma(z-2) \cdot \Omega((1 - \lambda)z) dz \\ & = \difrac{s^2}{2\pi i} \diint_{c - i\infty}^{c + i\infty} \difrac{s^{-z} \cdot \Gamma(z) \cdot \Omega((1 - \lambda)z)}{(z - 1) \cdot (z - 2)} dz. \end{split} \end{equation*}

For $\Re(z) > 0$, this function is clearly analytic. Now, for the points on the line $\Re(z) = 0$ (and points near the left of this line), we will use the information from theorem \ref{wellknowntheorem2}. Write $$\Upsilon(s, z) = \difrac{s^{-z} \cdot \Gamma(z) \cdot \Omega((1 - \lambda)z)}{(z - 1) \cdot (z - 2)}$$ to simplify the notation. By the residue theorem, we can write \begin{equation*} \begin{split} \diint_{c - i\infty}^{c + i\infty} \Upsilon(s, z) dz = 2\pi i [\Res(\Upsilon, 2) + \Res(\Upsilon, 1) + \Res(\Upsilon, 0)] + \diint_{\mathcal{C}} \Upsilon(s, z) dz, \end{split} \end{equation*} where $\mathcal{C}$ is the contour given by $$\mathcal{C} = \left\{\sigma(t) + it; \ t \in \R, \ \sigma(t) = \frac{-a}{2\log(2 + |t|)}\right\}.$$

Note that contour $\mathcal{C}$ lies within the set $\mathcal{V}$ of theorem \ref{wellknowntheorem2}. Calculating the residues, one has: 

\noindent $\bullet \ \Res(\Upsilon, 2) = \dilim_{z \to 2} (z - 2) \cdot \Upsilon(s, z) = \difrac{s^{-2} \cdot \Gamma(2) \cdot \Omega((1 - \lambda) \cdot 2)}{2 - 1} = \difrac{\Omega(2 \cdot (1 - \lambda))}{s^2}$

\noindent $\bullet \ \Res(\Upsilon, 1) = \dilim_{z \to 1} (z - 1) \cdot \Upsilon(s, z) = \difrac{s^{-1} \cdot \Gamma(1) \cdot \Omega((1 - \lambda) \cdot 1)}{1 - 2} = -\difrac{\Omega(1 - \lambda)}{s}$

\noindent $\bullet \ \Res(\Upsilon, 0) = \dilim_{z \to 0} z \cdot \Upsilon(s, z) = \difrac{s^0 \cdot \Omega((1 - \lambda) \cdot 0)}{(-1) \cdot (-2)} \cdot \dilim_{z \to 0} z \cdot \Gamma(z) = \difrac{\Omega(0)}{2}$

\noindent so that \begin{equation*} \begin{split} \mathscr{M}^{-1}[\Gamma(z) \Omega((1 - \lambda)(z + 2))] & = \difrac{s^2}{2\pi i} \diint_{c - i\infty}^{c + i\infty} \Upsilon(s, z) dz \\ & = \Omega(2 \cdot (1 - \lambda)) - s \cdot \Omega(1 - \lambda) + s^2 \cdot \difrac{\Omega(0)}{2} + \difrac{s^2}{2\pi i} \diint_{\mathcal{C}} \Upsilon(s, z) dz. \end{split} \end{equation*}

Thus, it suffices to show that the remaining integral is absolutely convergent for $s \in \mathcal{V}$. For each of these $s$, we have:

\noindent $\bullet \ |s^{-z}| = |s^{-\sigma(t) - it}| = |s|^{-\sigma(t)} \cdot \exp[t \cdot \arg(s)]$

\noindent $\bullet \ |\Gamma(z)| = |\Gamma(\sigma(t) + it)| \sim \sqrt{2\pi} \cdot |t|^{\sigma(t) - 1/2} \cdot \exp\left(-\difrac{\pi}{2} \cdot |t|\right)$ as $|t| \to \infty$ (uniformly for $\sigma$ in compact 

\noindent intervals, by Stirling's approximation)

\noindent $\bullet \ \Omega((1 - \lambda)z) \sim \log(|t|)$ \ \ \ \ \ (by theorem \ref{wellknowntheorem2})

\noindent $\bullet |(z - 1) \cdot (z - 2)| \sim t^2 + \difrac{5}{2}$

With all this in hand, we finally have \begin{equation*} \begin{split} |\Upsilon(s, z)| & = \difrac{|s^{-z}| \cdot |\Gamma(z)| \cdot |\Omega((1 - \lambda)z)|}{|(z - 1) \cdot (z - 2)|} \\ & \sim \difrac{\sqrt{2\pi} \cdot |s|^{-\sigma(t)} \cdot \exp\left[t \cdot \arg(s) - \difrac{\pi}{2} \cdot |t|\right] \cdot |t|^{-\sigma(t) - 1/2} \cdot \log(|t|)}{t^2 + 5/2} \\ & \xrightarrow[t \to \infty]{} \sqrt{2\pi} \cdot \exp\left(\difrac{a}{2}\right) \cdot \difrac{|t|^{-1/2} \cdot \log(|t|)}{t^2 + 5/2}, \end{split} \end{equation*} ensuring that the integral $$\diint_{-\infty}^{\infty} |\Upsilon(s, z)| dt \sim \diint_{-\infty}^{\infty} \difrac{|t|^{-1/2} \cdot \log(|t|)}{t^2 + 5/2} dt$$ converges. Furthermore, the same applies to the derivative of $\Upsilon(s, z)$ with respect to $s$: \begin{equation*} \begin{split} \Bigg|\difrac{d}{ds} \Upsilon(s, z)\Bigg| & = \Bigg| \difrac{d}{ds} \difrac{s^{-z} \cdot \Gamma(z) \cdot \Omega((1 - \lambda)z)}{(z - 1) \cdot (z - 2)} \Bigg| \\ & = \Bigg| \difrac{-z \cdot s^{- z - 1} \cdot \Gamma(z) \cdot \Omega((1 - \lambda)z)}{(z - 1) \cdot (z - 2)}\Bigg| \\ & = |-z| \cdot |s^{-1}| \cdot |\Upsilon(s, z)| \\ & \sim \difrac{|t|}{|s|} \cdot \sqrt{2\pi} \cdot \exp\left(\difrac{a}{2}\right) \cdot \difrac{|t|^{-1/2} \cdot \log(|t|)}{t^2 + 5/2} \\ & = \difrac{\sqrt{2\pi}}{|s|} \cdot \exp\left(\difrac{a}{2}\right) \cdot \difrac{|t|^{1/2} \cdot \log(|t|)}{t^2 + 5/2}, \end{split} \end{equation*} ensuring that the integral $$\diint_{-\infty}^{\infty} \Bigg|\difrac{d}{ds} \Upsilon(s, z)\Bigg| dt \sim \diint_{-\infty}^{\infty} \difrac{|t|^{1/2} \cdot \log(|t|)}{t^2 + 5/2} dt$$ also converges.

Since the integral $\diint_{\mathcal{C}} \Upsilon(s, z) dz$ and its derivative (with respect to $s$) converge absolutely and uniformly for $s$ in any compact subset of the half-plane $\Re(s) \geq 0$ (noting that the $s^2$ factor at the front compensates for the $s^{-1}$ term in the derivative as $s \to 0$), it follows from the differentiation theorem under the integration sign that the function$$s \mapsto \difrac{s^2}{2\pi i} \int_{\mathcal{C}} \Upsilon(s, z) dz $$is analytic in an open neighborhood of the half-plane $\Re(s) \geq 0$. \end{proof}

\end{document}
